\documentclass[12pt]{article}

% Packages for math, formatting, and layout
\usepackage{amsmath}
\usepackage{amssymb}
\usepackage{geometry}
\usepackage{hyperref}
\usepackage{enumitem}

% Set page margins
\geometry{a4paper, margin=1in}

\title{\textbf{Policy Gradient Methods}}
\author{}
\date{}

\begin{document}

\maketitle

\vspace{0.5cm}

\section*{Policy Gradient Methods}
\begin{itemize}[label=$\bullet$]
    \item Policy Gradient methods bypass value functions as the primary decision-maker.
    \item Instead, they directly parameterize the policy and optimize it using gradient ascent.
    \item In Policy Gradient (PG) methods, we represent the policy $\pi$ using a parameterized function (usually a neural network) with weights $\theta$.
    \item The policy outputs a probability distribution over actions for a given state:
    $$\pi_\theta(a \mid s) = \mathbb{P}[A_t = a \mid S_t = s, \theta]$$
\end{itemize}

\subsection*{Advantages over action-value methods like DQN:}
\begin{enumerate}
    \item DQN struggles with continuous actions because finding the maximum Q-value requires an exhaustive search or complex optimization at every single step.
    \begin{itemize}[label=$\circ$]
        \item PG methods simply output the parameters of a continuous distribution (like the mean and variance of a Gaussian).
    \end{itemize}
    \item Q-learning always learns a deterministic policy.
    \begin{itemize}[label=$\circ$]
        \item PG methods can learn optimal stochastic policies too.
    \end{itemize}
    \item Small changes in Q-values can cause the action selection to drastically flip, causing unstable training.
    \begin{itemize}[label=$\circ$]
        \item Small changes in policy parameters $\theta$ smoothly change the action probabilities.
    \end{itemize}
\end{enumerate}

\section*{The Policy Gradient Theorem (Derivation)}

\subsection*{1. Defining Aim}
\begin{itemize}[label=$\bullet$]
    \item We want our agent to maximize its total reward.
    \item Let's call a full episode a \textit{trajectory}, represented by $\tau$. 
    \item The total reward for a trajectory is $R(\tau)$.
    \item The probability of a specific trajectory happening, given our current policy weights $\theta$, is $P(\tau \mid \theta)$.
    \item Our objective function $J(\theta)$ is simply the expected return over all possible trajectories:
    $$J(\theta) = \int P(\tau \mid \theta) R(\tau) d\tau$$
\end{itemize}

\subsection*{2. Maximizing the Objective}
\begin{itemize}[label=$\bullet$]
    \item To maximize this objective using gradient ascent, we need to find the gradient (derivative) with respect to our policy parameters $\theta$.
    \item Since the reward $R(\tau)$ comes from the environment and does not depend on $\theta$, we only take the gradient of the probability:
    $$\nabla_\theta J(\theta) = \int \nabla_\theta P(\tau \mid \theta) R(\tau) d\tau$$
    \item Now we need this equation to be an expectation so we can estimate it by simply letting our agent play the game.
    \item We know the log-derivative trick:
    $$\nabla \log(x) = \frac{\nabla x}{x}$$
    \item On rearranging, we get:
    $$\nabla x = x \nabla \log(x)$$
    \item Now, replacing $x$ with our probability $P(\tau \mid \theta)$:
    $$\nabla_\theta P(\tau \mid \theta) = P(\tau \mid \theta) \nabla_\theta \log P(\tau \mid \theta)$$
    \item Substituting this back into our integral:
    $$\nabla_\theta J(\theta) = \int P(\tau \mid \theta) \nabla_\theta \log P(\tau \mid \theta) R(\tau) d\tau$$
    \item Because we have $P(\tau \mid \theta)$ multiplying everything again, this is now an expectation:
    $$\nabla_\theta J(\theta) = \mathbb{E}_{\tau \sim \pi_\theta} \left[ \nabla_\theta \log P(\tau \mid \theta) R(\tau) \right]$$
\end{itemize}

\subsection*{3. What $\nabla_\theta \log P(\tau \mid \theta)$ actually is}
\begin{itemize}[label=$\bullet$]
    \item What is the probability of a trajectory $P(\tau \mid \theta)$? It is the product of three things:
    \begin{enumerate}
        \item The probability of the starting state: $P(s_0)$
        \item The probability of the agent's actions (our policy): $\pi_\theta(a_t \mid s_t)$
        \item The probability of the environment's transitions: $P(s_{t+1} \mid s_t, a_t)$
    \end{enumerate}
    \item Therefore, on simplifying with logarithms and substituting:
    $$\log P(\tau \mid \theta) = \log P(s_0) + \sum_{t=0}^T \log \pi_\theta(a_t \mid s_t) + \sum_{t=0}^T \log P(s_{t+1} \mid s_t, a_t)$$
\end{itemize}

\subsection*{4. Eliminating the Environment Dynamics}
\begin{itemize}[label=$\bullet$]
    \item We take the gradient of the entire log sum with respect to $\theta$.
    \item The environment's starting state ($s_0$) and transition probabilities ($P(s_{t+1} \mid s_t, a_t)$) have absolutely nothing to do with our policy weights $\theta$. 
    \item Therefore, their derivatives are exactly zero!
    \item All that survives the gradient is our policy:
    $$\nabla_\theta \log P(\tau \mid \theta) = \sum_{t=0}^T \nabla_\theta \log \pi_\theta(a_t \mid s_t)$$
\end{itemize}

\subsection*{5. The Final Equation}
$$\nabla_\theta J(\theta) = \mathbb{E}_{\tau \sim \pi_\theta} \left[ \sum_{t=0}^T \nabla_\theta \log \pi_\theta(a_t \mid s_t) R(\tau) \right]$$

\begin{itemize}[label=$\bullet$]
    \item To improve your policy, you take the gradient of the log probability of your actions and multiply it by the total reward $R(\tau)$. 
    \item If the reward is high, the gradient steps in a direction that makes those actions more likely. 
    \item If the reward is low or negative, it pushes those actions away.
\end{itemize}

\section*{The REINFORCE Algorithm}
\begin{itemize}[label=$\bullet$]
    \item REINFORCE is the most fundamental Monte Carlo Policy Gradient algorithm.
    \item It takes the Policy Gradient Theorem and applies it in a practical, sample-based way.
    \item Instead of computing the true mathematical expectation, the agent plays an entire episode to generate a sample trajectory.
    \item At each time step $t$, it updates the policy weights $\theta$ using the actual return from that time step onward ($G_t$).
    \item $G_t$ acts as a scalar multiplier:
    \begin{itemize}[label=$\circ$]
        \item If the return $G_t$ was highly positive, we take a big step in that direction, making that action much more likely in the future.
        \item If the return $G_t$ was negative (or very small), we push the probability of that action down.
    \end{itemize}
\end{itemize}

\noindent\textbf{The REINFORCE Update Rule:}
$$\theta \leftarrow \theta + \alpha \nabla_\theta \log \pi_\theta(A_t \mid S_t) G_t$$

\section*{REINFORCE with Baseline and Variance Reduction}
\begin{itemize}[label=$\bullet$]
    \item The standard REINFORCE algorithm works, but it suffers from extreme variance.
    \item Because $G_t$ depends on a long sequence of random actions and state transitions, the returns can fluctuate wildly from one episode to the next.
    \item We can significantly reduce this variance by subtracting a \textbf{baseline} $b(s)$ from the return.
    \item The baseline can be any function, as long as it does not depend on the action $a$.
\end{itemize}

\noindent\textbf{The Baseline Update Rule:}
$$\theta \leftarrow \theta + \alpha \nabla_\theta \log \pi_\theta(A_t \mid S_t) (G_t - b(S_t))$$

\subsection*{Why the Baseline is Unbiased?}
\begin{itemize}[label=$\bullet$]
    \item Subtracting a random number from our return would not break the math and bias our gradient because the expected value of the baseline term is exactly zero.
\end{itemize}

\subsection*{Variance Reduction}
\begin{itemize}[label=$\bullet$]
    \item While any baseline works mathematically, the most effective baseline is the state-value function $V(s)$, which represents the average expected return from state $s$.
    \item When we substitute $V(s)$ into our equation, the multiplier becomes:
    $$A_t = G_t - V(S_t)$$
    This term is called the \textbf{Advantage}.
    \item If $G_t > V(S_t)$, the action we took resulted in a return better than average. The Advantage is positive, and the probability of that action is pushed heavily upwards.
    \item If $G_t < V(S_t)$, the action resulted in a return worse than average. The Advantage is negative, and the probability of that action is penalized and pushed down.
    \item By centering the returns around zero, the variance of the gradients shrinks dramatically, allowing the policy network to learn much faster and far more stably.
\end{itemize}

\end{document}